% Abstract for IMU: at most 250 words after macro expansion (checked by check_imu.py; also by plain whitespace count).
\noindent\textbf{Background:} Chest radiograph classifiers are often trained on pooled public collections in which disease class and image source are confounded, so we asked where source information resides in frozen features and whether it can be removed.

\noindent\textbf{Methods:} We rebuilt a six-class benchmark from official sources with grouped splits and an external cohort. On frozen DenseNet-121 features we measured source and disease information in six spatial-frequency bands, applied linear erasure (INLP and LEACE), and held out each source in turn with linear heads. Decision rules were fixed in advance and recorded in the repository.

\noindent\textbf{Results:} On the internal test split the model reached \RTestAcc\% accuracy, yet it called only \RExtSpecK/\RExtSpecN{} of the normal external scans Normal. Source balanced accuracy was \RSpecMainTestFullSrcBA\% (chance 16.7\%) and exceeded chance in every band, from \RSpecMainTestPassASrcBA\% to \RSpecMainTestPassFSrcBA\%; removing one band lowered it by at most \RSpecMainTestStopFDeltaS{} points. Twenty erasure iterations reduced it from \RInlpImgTestSrcBAKZero\% to \RInlpImgTestSrcBAKStar\%, not to chance, and also lowered disease scores. For four of five sources, the macro-F1 of a linear head trained on the other sources was \RLosoImgCohenGap{} to \RLosoImgTBXkGap{} points below the in-distribution reference. A band-stop retraining was not run, because the go/no-go rule fixed in advance returned \RBsDecision.

\noindent\textbf{Conclusions:} In these frozen features, source information is spread over all frequency bands and entangled with disease information, so neither band removal nor linear erasure is a cheap remedy. To our knowledge, within the searches described, this band-by-band localisation has not been reported for chest radiographs. Evidence that the model uses the source is indirect.
